\usepackage[
    paperwidth=8.5in,
    paperheight=11in,
    margin=1.25in,
    marginparwidth=0.75in,
    marginparsep=0.1in,
    footskip=0.25in
]{geometry}

\renewcommand{\baselinestretch}{1.}
\setlength{\parindent}{2em}
\setlength{\parskip}{2pt}
\usepackage[shortlabels]{enumitem}
\setlist{
    itemsep=1ex,
    listparindent=\parindent,
    parsep=0em,
    topsep=2ex,
}
\setcounter{tocdepth}{1}
\allowdisplaybreaks[1]

\usepackage{xpatch}
\usepackage{amsfonts}
\usepackage{amsmath}
\usepackage{amssymb}
\usepackage{amsthm,thmtools,thm-restate}
\usepackage{bm}
\usepackage{mathrsfs}
\usepackage{mathtools}

\usepackage{booktabs}
\usepackage{comment}
\usepackage{float}
\usepackage{graphicx}
\usepackage{multirow}
\usepackage{tikz,pgfplots}
\pgfplotsset{compat=1.18}
\usepgfplotslibrary{groupplots}
\usetikzlibrary{calc}

\usepackage{silence}
\WarningFilter[pdftoc]{hyperref}{Token not allowed in a PDF string}

\usepackage{caption}
\captionsetup[figure]{skip=10pt, width=12.25cm}
\captionsetup{width=11cm}

\usepackage[
    backend=biber,
    style=numeric,
    giveninits=true,
    sorting=nyt,
    maxnames=100,
    backref=false,
    doi=false,
    url=false,
    isbn=false,
]{biblatex}
\addbibresource{ref.bib}
\renewcommand*{\bibfont}{\footnotesize}
\renewbibmacro{in:}{}
\DeclareFieldFormat[
	article,
	inbook,
	incollection,
	inproceedings,
	patent,
	thesis,
	unpublished
]{title}{#1}

\newcommand\linkcolor{black}

\usepackage{hyperref}
\hypersetup{
    colorlinks,
    linktoc=all,
    citecolor=black,
    linkcolor=\linkcolor,
    urlcolor=\linkcolor
}
\usepackage[capitalize,nameinlink,noabbrev]{cleveref}
\crefname{claim}{Claim}{Claims}
\crefname{fact}{Fact}{Facts}
\crefname{inequality}{Inequality}{Inequalities}
\creflabelformat{inequality}{#2(#1)#3}
\crefname{lemma}{Lemma}{Lemmas}
\crefname{subsection}{Subsection}{Subsections}
\newcommand{\creflastconjunction}{, and\nobreakspace}

\makeatletter
\def\cref@short@range@type{}
\def\cref@lemma@type{lemma}
\def\cref@section@type{section}
\newcommand{\cref@reinsert@short@range}[1]{%
  \cref@isstackfull{\cref@consecutive@stack}%
  \@whilesw\if@cref@stackfull\fi{%
    \edef\@tempa{\cref@stack@top{\cref@consecutive@stack}}%
    \expandafter\cref@stack@push\expandafter{\@tempa}{#1}%
    \cref@stack@pop{\cref@consecutive@stack}%
    \cref@isstackfull{\cref@consecutive@stack}}}
\xpatchcmd{\cref@processconsecutive}
  {\let#4\relax%
  #5=1\relax}
  {\let#4\relax%
  #5=1\relax%
  \cref@stack@init{\cref@consecutive@stack}}
  {}
  {\PackageWarning{preamble}{Could not initialize cleveref range stack}}
\xpatchcmd{\cref@processconsecutive}
  {\advance#5 1\relax%
      \let#4\@nextref%
      \cref@stack@pop{#2}}
  {\advance#5 1\relax%
      \let#4\@nextref%
      \expandafter\cref@stack@push\expandafter{\@nextref}{\cref@consecutive@stack}%
      \cref@stack@pop{#2}}
  {}
  {\PackageWarning{preamble}{Could not record cleveref range members}}
\xpatchcmd{\@cref}
  {\ifnum\count@consecutive=2\relax%
        \expandafter\cref@stack@push\expandafter{\@endref,}{\@refsubstack}%
        \let\@endref\relax%
        \count@consecutive=1\relax%
      \fi}
  {\ifnum\count@consecutive>1\relax%
        \ifnum\count@consecutive=2\relax%
          \cref@reinsert@short@range{\@refsubstack}%
          \let\@endref\relax%
          \count@consecutive=1\relax%
        \else%
          \ifnum\count@consecutive=3\relax%
            \expandafter\cref@gettype\expandafter{\@beginref}{\cref@short@range@type}%
            \ifx\cref@short@range@type\cref@lemma@type%
              \cref@reinsert@short@range{\@refsubstack}%
              \let\@endref\relax%
              \count@consecutive=1\relax%
            \else%
              \ifx\cref@short@range@type\cref@section@type%
                \cref@reinsert@short@range{\@refsubstack}%
                \let\@endref\relax%
                \count@consecutive=1\relax%
              \fi%
            \fi%
          \fi%
        \fi%
      \fi}
  {}
  {\PackageWarning{preamble}{Could not adjust cleveref range threshold}}
\makeatother

\renewcommand\qedsymbol{$\blacksquare$}
\theoremstyle{plain}
\newtheorem{theorem}{Theorem}[section]
\newtheorem{claim}[theorem]{Claim}
\newtheorem{corollary}{Corollary}[theorem]
\newtheorem{fact}[theorem]{Fact}
\newtheorem{lemma}[theorem]{Lemma}
\newtheorem{proposition}[theorem]{Proposition}
\newtheorem{example}[theorem]{Example}
\newtheorem{definition}[theorem]{Definition}

\theoremstyle{definition}
\newtheorem{conjecture}[theorem]{Conjecture}
\newtheorem{remark}[theorem]{Remark}

\newcommand\E{{\mathbb E}}
\newcommand\N{{\mathbb N}}
\newcommand\bbP{{\mathbb P}}
\newcommand\R{{\mathbb R}}
\newcommand\Z{{\mathbb Z}}
\newcommand\bsone{\boldsymbol{1}}
\newcommand\cB{{\mathcal B}}
\newcommand\cC{{\mathcal C}}
\newcommand\cE{{\mathcal E}}
\newcommand\cF{{\mathcal F}}
\newcommand\cG{{\mathcal G}}
\newcommand\cH{{\mathcal H}}
\newcommand\cI{{\mathcal I}}
\newcommand\cJ{{\mathcal J}}
\newcommand\cK{{\mathcal K}}
\newcommand\cM{{\mathcal M}}
\newcommand\cN{{\mathcal N}}
\newcommand\cP{{\mathcal P}}
\newcommand\cR{{\mathcal R}}
\newcommand\cS{{\mathcal S}}
\newcommand\cT{{\mathcal T}}
\newcommand\cV{{\mathcal V}}
\newcommand\cW{{\mathcal W}}
\newcommand\cX{{\mathcal X}}
\newcommand\cY{{\mathcal Y}}
\newcommand\cZ{{\mathcal Z}}
\newcommand\frp{{\mathfrak p}}
\newcommand\frP{{\mathfrak P}}
\newcommand\frt{{\mathfrak t}}
\newcommand\sB{{\mathscr B}}
\newcommand\sD{{\mathscr D}}
\newcommand\sE{{\mathscr E}}
\newcommand\sP{{\mathscr P}}
\newcommand\sR{{\mathscr R}}
\newcommand{\red}{{\normalfont\texttt r}}
\newcommand{\green}{{\normalfont\texttt g}}
\newcommand{\blue}{{\normalfont\texttt b}}

\newcommand\erdos{Erd\H{o}s}
\newcommand\renyi{R\'enyi}
\newcommand\erdosrenyi{\erdos--\renyi}

\makeatletter
\newcommand{\nHookrightarrow}{\mathrel{\mathpalette\@nHookrightarrow\relax}}
\newcommand{\@nHookrightarrow}[2]{%
  \begin{tikzpicture}[baseline={(H.base)}]
    \node[inner sep=0pt, outer sep=0pt] (H) {$#1\hookrightarrow$};
    \draw[line cap=round, line width=.07em]
      ([xshift=.425em,yshift=.15ex]H.south west) --
      ([xshift=-.525em,yshift=-.055ex]H.north east);
  \end{tikzpicture}%
}
\makeatother

\pgfmathsetmacro{\opacity}{50}
\colorlet{lightred}{red!\opacity}
\colorlet{lightgreen}{green!\opacity}
\definecolor{lightblue}{RGB}{144,202,249}

\newcommand{\grg}{\tikz[baseline=-.6ex,x=1em,y=1em]{%
  \draw[lightred, line width=3pt,line cap=round] (0,0)--(1.25,0);
  \fill[lightgreen] (0,0) circle (3pt);
  \fill[lightgreen] (1.25,0) circle (3pt);
}}

\newcommand{\gbg}{\tikz[baseline=-.6ex,x=1em,y=1em]{%
  \draw[lightblue, line width=3pt,line cap=round] (0,0)--(1.25,0);
  \fill[lightgreen] (0,0) circle (3pt);
  \fill[lightgreen] (1.25,0) circle (3pt);
}}

\newcommand{\brb}{\tikz[baseline=-.6ex,x=1em,y=1em]{%
  \draw[lightred, line width=3pt,line cap=round] (0,0)--(1.25,0);
  \fill[lightblue] (0,0) circle (3pt);
  \fill[lightblue] (1.25,0) circle (3pt);
}}

\newcommand{\rrbtriangle}{\tikz[baseline=-0.15ex,x=1em,y=1em,scale=0.75]{%
  \coordinate (A) at (0,0);
  \coordinate (B) at (1.25,0);
  \coordinate (C) at (.625,1.08);
  \draw[lightblue,line width=2.5pt,line cap=round] (A)--(B); % AB
  \draw[lightred,line width=2.5pt,line cap=round] (B)--(C); % BC
  \draw[lightred,line width=2.5pt,line cap=round] (C)--(A); % CA

  \fill[lightblue] (C) circle (3pt);
  \filldraw[fill=white, draw=black, very thin] (A) circle (3pt);
  \filldraw[fill=white, draw=black, very thin] (B) circle (3pt);
}}

\newcommand{\rrgtriangle}{\tikz[baseline=-0.15ex,x=1em,y=1em,scale=0.75]{%
  \coordinate (A) at (0,0);
  \coordinate (B) at (1.25,0);
  \coordinate (C) at (.625,1.08);
  \draw[lightgreen,line width=2.5pt,line cap=round] (A)--(B); % AB
  \draw[lightred,line width=2.5pt,line cap=round] (B)--(C); % BC
  \draw[lightred,line width=2.5pt,line cap=round] (C)--(A); % CA

  \fill[lightgreen] (C) circle (3pt);
  \filldraw[fill=white, draw=black, very thin] (A) circle (3pt);
  \filldraw[fill=white, draw=black, very thin] (B) circle (3pt);
}}

\newcommand{\exampletriangle}{\tikz[baseline=-0.15ex,x=1em,y=1em,scale=0.75]{%
  \coordinate (A) at (0,0);
  \coordinate (B) at (1.25,0);
  \coordinate (C) at (.625,1.08);
  \draw[lightred,line width=2.5pt,line cap=round] (A)--(B); % AB
  \draw[lightred,line width=2.5pt,line cap=round] (B)--(C); % BC
  \draw[lightred,line width=2.5pt,line cap=round] (C)--(A); % CA

  \fill[lightgreen] (A) circle (3pt);
  \fill[lightgreen] (B) circle (3pt);
  \fill[lightblue] (C) circle (3pt);
}}

\makeatletter
\newcommand{\vast}{\bBigg@{3}}
\makeatother

\newcommand\blambda{{\boldsymbol{\lambda}}}
\newcommand\bLambda{{\boldsymbol{\Lambda}}}
\newcommand\bsa{\boldsymbol{a}}
\newcommand\bsm{\boldsymbol{m}}
\newcommand{\deltatilde}{\tilde{\delta}}

\newcommand\close{{\hspace{0.2mm}\ensuremath{\mathsf{close}}}}
\newcommand\far{{\hspace{0.2mm}\ensuremath{\mathsf{far}}}}
\newcommand\high{{\hspace{0.2mm}\ensuremath{\mathsf{hi}}}}
\newcommand\mat{{\hspace{0.2mm}\ensuremath{\mathsf{mat}}}}
\newcommand\ind{{\hspace{0.2mm}\ensuremath{\mathrm{ind}}}}
\newcommand\spl{\ensuremath{\mathrm{spl}}}
\newcommand\degenerate{{\hspace{0.2mm}\ensuremath\mathsf{deg}}}
\newcommand\med{{\hspace{0.2mm}\ensuremath\mathsf{med}}}
\newcommand\res{{\hspace{0.2mm}\ensuremath\mathsf{res}}}
\newcommand\prof{{\hspace{0.2mm}\ensuremath\mathsf{prof}}}
\newcommand\nar{{\ensuremath\mathsf{nar}}}
\renewcommand\root{{\ensuremath\mathsf{root}}}
\newcommand\leftsf{{\hspace{0.2mm}\ensuremath\mathsf{left}}}
\newcommand\insf{{\hspace{0.2mm}\ensuremath\mathsf{in}}}
\newcommand\outsf{{\hspace{0.2mm}\ensuremath\mathsf{out}}}
\newcommand\Tail{{\ensuremath\mathsf{Tail}}}
\newcommand\Upper{{\ensuremath\mathsf{U}}}
\newcommand\Lower{{\ensuremath\mathsf{L}}}
\newcommand\bin{{\ensuremath\mathsf{bin}}}
\newcommand\Ent{{\ensuremath\mathsf{Ent}}}
\newcommand\Err{{\ensuremath\mathsf{Err}}}
\newcommand\ret{{\ensuremath\mathsf{ret}}}
\newcommand\matsf{{\ensuremath\mathsf{mat}}}
\newcommand\lgsf{{\ensuremath\mathsf{lg}}}
\newcommand\smsf{{\ensuremath\mathsf{sm}}}
\newcommand\sparse{{\ensuremath\mathsf{sp}}}

\DeclareMathOperator{\argmax}{argmax}
\DeclareMathOperator{\image}{image}
\DeclareMathOperator{\Bin}{Bin}
\DeclareMathOperator{\rand}{rand}

\DeclarePairedDelimiter\ceil{\lceil}{\rceil}
\DeclarePairedDelimiter\floor{\lfloor}{\rfloor}
\DeclarePairedDelimiter\abs{\lvert}{\rvert}
\DeclarePairedDelimiter\norm{\lVert}{\rVert}
\renewcommand\geq{\geqslant}
\renewcommand\leq{\leqslant}

\newcommand{\clean}{{\mathrm{cl}}}
\newcommand{\type}{{\mathrm{ty}}}
